% SPDX-License-Identifier: CC-BY-4.0
% Repository: https://github.com/zx1158/etam-model
% Version: v2.0 (2026-02-22)
% Language: English (中文注释已移除以确保编译)
% For GitHub rendering: see README.md

\documentclass[11pt]{article}
\usepackage[a4paper, left=2.5cm, right=2.5cm, top=2.5cm, bottom=2.5cm]{geometry}
\usepackage[utf8]{inputenc}      % 支持UTF-8编码
\usepackage[T1]{fontenc}         % 支持欧洲字符
\usepackage{times}               % Times字体
\usepackage{amsmath, amssymb, mathtools}
\usepackage{bm}
\usepackage{enumitem}
\usepackage{hyperref}
\usepackage{cleveref}
\usepackage{siunitx}
\usepackage{titlesec}
\usepackage{natbib}
\usepackage{graphicx}            % 添加图形支持
\usepackage{booktabs}            % 更好的表格
\usepackage{microtype}           % 微调排版

% 设置超链接
\hypersetup{
    colorlinks=true,
    linkcolor=blue,
    citecolor=blue,
    urlcolor=blue,
    pdftitle={ETAMv2.0: Energetic--Temporal--Environmental Activity Model},
    pdfauthor={ETAM Specification Consortium}
}

\sisetup{locale=US}

% Math operators
\DeclareMathOperator{\Var}{Var}
\DeclareMathOperator{\IE}{IE}
\DeclareMathOperator{\solv}{solv}
\DeclareMathOperator{\logE}{log}

% Symbols
\newcommand{\R}{\mathbb{R}}
\newcommand{\Z}{\mathbb{Z}}
\newcommand{\C}{\mathbb{C}}
\newcommand{\RePart}{\Re}
\newcommand{\ImPart}{\Im}
\newcommand{\eps}{\varepsilon}
\newcommand{\abs}[1]{\left|#1\right|}
\newcommand{\norm}[1]{\left\|#1\right\|}
\newcommand{\avg}[1]{\langle #1 \rangle}

% Bilingual terminology macros (英文版本)
\newcommand{\ATE}{\textbf{A--T--E} (Spatial-Temporal-Energetic)}
\newcommand{\XR}{\textbf{XR} (ATE Basis Encoding)}
\newcommand{\icR}{\textbf{icR} (Complex Generator: $f(R)=x(R)+i\,c(R)$)}

% 标题格式
\title{ETAMv2.0: Energetic--Temporal--Environmental Activity Model, Version 2.0}
\subtitle{Formal Specification Document (v2.0) — Not a derivation notebook nor implementation guide}
\author{ETAM Specification Consortium}
\date{2026-02-22}

\begin{document}

\maketitle

\section{Notation and Conventions}
\label{sec:notation}
All symbols are defined as follows unless otherwise noted:
\begin{itemize}[leftmargin=*, nosep]
    \item \ATE: primitive ontological basis — \textbf{A} for \emph{spatial-architectural activity}, \textbf{T} for \emph{temporal-slot activation}, \textbf{E} for \emph{energo-material response}
    \item \XR: discrete coordinate representation of the \ATE basis; $R = [\mathbf{n}_1,\dots,\mathbf{n}_N]^\top \in \Z^{N \times 3}$
    \item \icR: complex-valued generator $f(R) = x(R) + i\,c(R)$, where $x$ is the \emph{covariant observable}, $c$ the \emph{contravariant lifting potential}
\end{itemize}

\section{Introduction}
\label{sec:introduction}
The \textbf{Energetic--Temporal--Environmental Activity Model (ETAM)} is a physics-informed, differentiable framework for predicting macroscopic scalar outputs $Y$ from discrete-continuous hybrid inputs. Version 2.0 introduces unified slot-based encoding, complex-valued response functions, and adaptive multilevel hierarchy driven by the imaginary component.

\textbf{Clarification of novelty}: The \ATE\ triad is not a re-labeling of conventional physical quantities. It is a \emph{primitive ontological basis} defined by its role in the \icR\ generator: \textbf{A} encodes spatial-architectural constraints, \textbf{T} encodes temporal-slot resolution, \textbf{E} encodes energo-material coupling. No existing model uses this triad as a unified coordinate system.

This aligns with Rosen's relational ontology: \ATE\ does not describe \emph{what things are}, but \emph{how they co-constitute observables through constrained coordination} \cite{rosen1978}.

\section{Core Output Structure}
\label{sec:core-output}
The model defines a dual-domain mapping:
\begin{align}
    Y &= A \cdot T \cdot E \label{eq:output-physical-observable} \\
    f(R) &= x(R) + \mathrm{i}\, c(R) \label{eq:response-complex-generator}
\end{align}
where $Y \in \R_{\geq 0}$ is the physical observable (e.g., binding energy in \si{\electronvolt}), and $f(R) \in \C$ is the \icR\ generator. Alignment requires $\RePart[f(R)] \approx Y$, while $\ImPart[f(R)] = c(R) \geq 0$ governs scale-level selection.

\section{Input Representation $R$}
\label{sec:input-representation}
Each active unit $j$ (atom or functional group) is encoded as a discrete triple:
\begin{equation}
    \mathbf{n}_j = (n_{A,j},\, n_{T,j},\, n_{E,j}) \in \Z_{\geq 0}^3.
    \label{eq:unit-encoding}
\end{equation}
For $N$ units, the input tensor is:
\begin{equation}
    R = [\mathbf{n}_1,\, \mathbf{n}_2,\, \dots,\, \mathbf{n}_N]^\top \in \Z^{N \times 3}.
    \label{eq:input-tensor}
\end{equation}
$R$ is not a feature vector — it is the \emph{coordinate representation of the \ATE\ basis in discrete state space}. Each row $\mathbf{n}_j$ is a point in the \ATE\ lattice, where $n_{A,j}$, $n_{T,j}$, $n_{E,j}$ are \emph{mutually constraining coordinates}, not independent scalars.

\section{Component Definitions}
\label{sec:components}

\subsection{A-value: Spatial-Architectural Activity}
\label{subsec:a-value}
Let $\omega(p) \in \{0,1,2\}$ denote occupancy (empty/single electron/lone pair) at orbital slot $p$. Define slot weights $k(p)$:
\begin{equation}
    k(p) =
    \begin{cases}
        1 & \text{for } s\text{-orbitals}, \\
        2 & \text{for } p\text{-orbitals}, \\
        3 & \text{for } d\text{-orbitals}, \\
        4 & \text{for } f\text{-orbitals}.
    \end{cases}
    \label{eq:orbital-weights}
\end{equation}
The activity constant for unit $j$ is:
\begin{equation}
    C_{A,j} = \sum_{p \in \{s,p,d,f\}} \omega(p)\, k(p).
    \label{eq:ca-j}
\end{equation}
Then:
\begin{equation}
    A_j = \alpha_A \cdot C_{A,j}, \quad
    A = \frac{1}{N}\sum_{j=1}^N A_j,
    \label{eq:a-j-and-mean}
\end{equation}
with scaling factor $\alpha_A > 0$.

\subsection{T-value: Temporal-Slot Activation}
\label{subsec:t-value}
Six discrete time slots $k=0,\dots,5$ cover $\SI{1e-15}{\second}$ to $\SI{1e6}{\second}$. Let $\nu(k) \in \{0,1,2\}$ indicate activation level. Then:
\begin{equation}
    C_{T,j} = \sum_{k=0}^{5} \nu(k)\, w(k), \quad
    T_j = \alpha_T \cdot C_{T,j}, \quad
    T = \frac{1}{N}\sum_{j=1}^N T_j.
    \label{eq:t-j-and-mean}
\end{equation}

\subsection{E-value: Energo-Material Response Field}
\label{subsec:e-value}
Four environmental response slots: ionization ($e_I$), affinity ($e_A$), polarization ($e_P$), solvation ($e_S$). Let $\mu(e) \in \{0,1,2\}$. Define:
\begin{equation}
    C_{E,j} = \sum_{e \in \{I,A,P,S\}} \mu(e)\, w_e.
    \label{eq:ce-j}
\end{equation}
Environment-modulated field:
\begin{align}
    \Phi_I(\Xi) &= \frac{F^2}{\IE}, 
    \label{eq:phi-ionization} \\
    \Phi_P(\Xi) &= \frac{\epsilon - 1}{\epsilon + 2} F^2, 
    \label{eq:phi-polarization} \\
    \Phi_S(\Xi) &= \frac{\Delta G_{\solv}}{RT},
    \label{eq:phi-solvation}
\end{align}
and:
\begin{equation}
    E_j = \alpha_E \cdot C_{E,j} \cdot \Bigg[1 + \sum_{e} \underbrace{\frac{\mu(e)\, w_e}{C_{E,j}}}_{\beta_e(C_{E,j})} \cdot \Phi_e(\Xi) \Bigg], \quad
    E = \frac{1}{N}\sum_{j=1}^N E_j.
    \label{eq:e-j-and-mean}
\end{equation}

\section{Response Function $f(R)$}
\label{sec:response-function}

The \icR\ framework ($f(R) = x(R) + i\,c(R)$) is not an engineering choice for numerical stability. It is a \emph{structural necessity}: only complex-valued representation permits simultaneous encoding of covariant observables ($x$) and contravariant scale selectors ($c$) within a single differentiable map.

The icR mapping realizes a \emph{Connes-type spectral geometry} \cite{connes1994}: discrete \ATE\ coordinates $R$ define a finite spectral triple, whose complex extension $f(R)$ yields a differentiable continuum limit.

\subsection{Real Part (Covariant Observable)}
\label{subsec:real-part}
\begin{equation}
    x(R) = \alpha_x \cdot \avg{C_A} \cdot \avg{C_T} \cdot \avg{C_E},
    \label{eq:x-r-product-form}
\end{equation}
or alternatively via MLP:
\begin{equation}
    x(R) = \mathrm{MLP}_{\theta_x}\big( \logE(C_A+\eps),\, \logE(C_T+\eps),\, \logE(C_E+\eps) \big).
    \label{eq:x-r-mlp-form}
\end{equation}

\subsection{Imaginary Part (Contravariant Lifting Potential)}
\label{subsec:imag-part}
\begin{equation}
    c(R) = \lambda \cdot \Bigg[ \frac{\Var(C_A)}{\avg{C_A}^2} + \frac{\Var(C_T)}{\avg{C_T}^2} \Bigg] \cdot \big(1 + \eta \cdot S_{\mathrm{env}} \big),
    \label{eq:c-r-lifting}
\end{equation}
where $S_{\mathrm{env}} = \sum_e \Phi_e(\Xi)$. The non-negative imaginary part $c(R)$ implements Penrose's 'ladder of infinity' on the complex plane \cite{penrose2004}: each threshold $\tau_k$ marks a rung where the system's effective description ascends to a new ontological level.
\begin{itemize}[leftmargin=*, nosep]
    \item $c(R) < \tau_1$: Level 1 (molecular determinism)
    \item $\tau_1 \leq c(R) < \tau_2$: Level 2 (mesoscale clusters)
    \item $c(R) \geq \tau_2$: Level 3 (condensed-phase statistical fields)
\end{itemize}

\section{Training Objective}
\label{sec:training-objective}
\begin{equation}
    \mathcal{L} = \norm{ \RePart[f(R)] - Y }^2 + \gamma \cdot \mathcal{L}_{\mathrm{reg}}\big(c(R)\big),
    \label{eq:training-loss}
\end{equation}
where $\mathcal{L}_{\mathrm{reg}}$ enforces level consistency (e.g., entropy penalty on $c(R)$) or other domain-informed constraints.

\section{Constrained Inverse Queries}
\label{sec:inverse-queries}
The model supports \emph{constrained inverse queries}: given a partial configuration $\tilde{R}$ and target scalar $y_{\mathrm{target}}$, the set of completions $R_{\mathrm{miss}}$ satisfying both energy fidelity and level consistency is non-empty and computable under standard differentiable optimization.

This inverse capability is not an add-on feature. It is a \emph{logical consequence} of the \ATE--\icR\ duality: if $f$ is bijective over its domain (up to level constraints), then generation is the algebraic inverse of prediction \cite{wigner1960}.

\section{Version Identification}
\label{sec:version}
\textbf{ETAMv2.0} is formally defined by this specification (v2.0, 2026-02-22). Key normative innovations:
\begin{enumerate}[leftmargin=*, nosep]
    \item Unified discrete slot encoding across \textbf{A}, \textbf{T}, and \textbf{E} as a primitive ontological basis;
    \item Complex-valued separation of covariant observable ($x$) and contravariant lifting potential ($c$);
    \item Imaginary part-driven adaptive multilevel hierarchy grounded in Penrose's ladder;
    \item Verifiable scientific discovery via Wigner-style bidirectional invertibility.
\end{enumerate}

% 如果没有bib文件，可以注释掉以下两行
\bibliographystyle{unsrtnat}
\bibliography{etam-philosophy}

% 如果没有bib文件，使用以下替代：
%\begin{thebibliography}{9}
%\bibitem{rosen1978} Rosen, R. (1978). Fundamentals of Measurement and Representation of Natural Systems.
%\bibitem{connes1994} Connes, A. (1994). Noncommutative Geometry.
%\bibitem{penrose2004} Penrose, R. (2004). The Road to Reality: A Complete Guide to the Laws of the Universe.
%\bibitem{wigner1960} Wigner, E. P. (1960). The Unreasonable Effectiveness of Mathematics in the Natural Sciences.
%\end{thebibliography}

\end{document}
